\documentclass[11pt]{article}
\usepackage[a4paper,margin=1in]{geometry}
\usepackage{amsmath,amssymb}
\usepackage{booktabs}
\usepackage{hyperref}
\usepackage{enumitem}
\usepackage{microtype}
\usepackage[numbers]{natbib}

\title{Conflux: Principal-Context Security, Evaluation and Verification for LLM Agents\\\large Work-in-progress fourth-year paper}
\author{Raya Buckley}
\date{29 July 2026}

\begin{document}
\maketitle

\begin{abstract}
Large-language-model agents combine information controlled by different
principals while performing externally visible actions. Previous work
introduced Influence Tracking with Extrapolated Security (ITES) and the
System-Level Evaluator for Defences (SLED). Conflux develops those ideas into a
provider- and benchmark-independent research framework. The current repository
has one canonical, fail-closed security kernel with immutable provenance,
independent read, authorisation, visibility and consent decisions, versioned
decision evidence, and bounded explicit-state verification. Architectural
repair is tested, but real-model, external-benchmark and solver-backed claims
remain gated on reproducible result artefacts. This paper distinguishes
archived evidence, current implementation evidence and planned experiments;
numerical placeholders are filled only by generated artefacts.
\end{abstract}

\section{Introduction}
LLM agents process requests, retrieved documents and tool output controlled by
principals with different authority. An external author should not acquire the
initiating Principal's permissions merely by influencing a model context.
Model-level prompt-injection defences improve instruction following, while
system-level approaches constrain effects externally
\citep{camel,designpatterns,struq,spotlighting}.

For Principal Context $PC$ and action $a$, the restricted ITES rule is
\[
  \operatorname{Allow}(PC,a) \Longleftrightarrow
  PC \neq \varnothing \land \forall p \in PC,\;P(p,a).
\]
The explicit non-empty condition prevents vacuous authority. Provenance grows
monotonically, while organisational authority is supplied by an independent
policy decision point.

\section{Previous-year contribution and archived evidence}
The previous prototype accumulated authorship provenance and intersected
permissions. Its bounded SLED results apply to that implementation and model.
The archived source and PDF remain under \texttt{paper/} with checksums. They
are not evidence for the refactored kernel. Current experiments use separately
named legacy-reproduction and canonical suites.

\section{Related work and positioning}
CaMeL isolates control and data flows, while privilege-oriented systems such as
Progent mediate tool calls with symbolic policy \citep{camel,progent}. PACT
tracks argument-level provenance and FORGE enforces history-sensitive Datalog
policies \citep{pact,forge}. AgentDojo supplies realistic prompt-injection tasks
and metrics \citep{agentdojo}. Conflux therefore does not claim priority for
system-level mediation, provenance or privilege control. Its narrower research
focus is collective Principal Context derived from authenticated provenance,
interpreted through organisational policy, and evaluated with a shared
operational and verification transition model.

\section{Canonical architecture}
Immutable domain values represent Principals, Principal Contexts, permissions,
resources, artefacts, provenance, sessions, actions and environment snapshots.
Policy ports decide read access, authorisation, visibility and consent
independently. The ITES kernel is pure: model proposals cannot assert or narrow
their Principal Context, alternative branches share an immutable parent, and
effect execution requires the exact decision certificate.

\subsection{Proposal semantics}
A versioned proposal batch declares either \texttt{ALTERNATIVES} or
\texttt{ORDERED\_PLAN}. Alternatives are independently explored from one
parent. Ordered steps propagate state but are re-authorised immediately before
execution; no batch grants future authority.

\subsection{Evidence and failure semantics}
Reports distinguish proposed, authorised, blocked, executed, provider-failed
and incomplete actions. Policy errors, malformed proposals, missing provenance
and exhausted bounds are explicit fail-closed outcomes. Securely blocked
attacks are diagnostics, not executed invariant failures.

\section{Evaluation methodology}
The deterministic semantic corpus covers empty and mixed contexts,
author/reader inversions, derivation, nested accumulation, branch isolation,
ordered plans, consent, visibility, unsupported delegation, policy changes and
provider failure. Executable defective variants must yield minimal
counterexamples.

Native SLED uses versioned legacy-reproduction and corrected canonical suites.
SLED-MC performs breadth-first exploration with canonical state keys,
deduplication, predecessor edges and explicit \texttt{SAFE},
\texttt{BOUNDED\_SAFE}, \texttt{UNSAFE} and \texttt{UNKNOWN} verdicts.

Real-model manifests record model revision, endpoint configuration, sampling,
structured-output mode, raw-response hashes, token use and latency. AgentDojo
integration pins an upstream revision, preserves upstream IDs, and reports
native and Conflux metrics without hiding added provenance assumptions.

\section{Formal verification and planning}
SLED-V accepts a finite serialisable transition IR rather than arbitrary
Python. Initial properties cover unauthorised effects and reads, provenance
monotonicity, absence of implicit delegation and bounded resource use. Bounded
SMT results retain solver artefacts; unbounded claims require an
IC3/PDR-capable backend and runtime-to-IR conformance.

Authority-constrained planning represents typed operations, preconditions,
outcomes, required reads and permissions. Security is a hard constraint.
Selection may then minimise authority footprint, sensitive observations,
calls, latency and irreversible effects.

\section{Results}
\textbf{Repository validation:} \emph{generated result pending}.\\
\textbf{Semantic conformance:} \emph{generated result pending}.\\
\textbf{Native SLED:} \emph{generated result pending}.\\
\textbf{SLED-MC/SLED-V:} \emph{generated result pending}.\\
\textbf{Real models and AgentDojo:} \emph{generated result pending}.

No numerical claim is entered manually. Tables and figures are generated from
retained, versioned result JSON and include input checksums.

\section{Limitations}
Conflux assumes sound authentication, provenance, policy state and complete
mediation. It does not prevent actions already authorised for every influencing
Principal and does not guarantee subjective intent alignment. Conservative
contexts can over-approximate influence. Delegation remains unsupported.
Authorised-read safety is not noninterference. Formal results are scoped to
their model, backend and conformance evidence. Provider adapters are research
prototypes, not production isolation boundaries.

\section{Conclusion}
Conflux now has one tested security semantics and a bounded state-exploration
foundation. The remaining contribution is evidence: reproducible native,
real-model, external-benchmark and formal-verification artefacts whose claims
remain tied to explicit assumptions and bounds.

\bibliographystyle{plainnat}
\bibliography{references}
\end{document}
